\documentclass[10pt,a4paper]{amsart}
\usepackage[margin=19mm]{geometry}
\usepackage{amsmath,amssymb,mathtools}
\usepackage[scr=rsfs,cal=euler]{mathalfa}
\usepackage{xfrac}
\usepackage[unicode,hidelinks,bookmarks=false]{hyperref}
\newcommand{\ie}{i.e.\ }
\newcommand{\cf}{cf.\ }
\newcommand{\resp}{respectively\ }
\newcommand{\Subsection}{\subsection}
\providecommand{\nobreakdash}{\nobreak\hbox{-}}
\setlength{\emergencystretch}{2em}
\multlinegap0pt
\usepackage{amssymb}
\usepackage{enumerate}
\usepackage{xcolor}
\usepackage{tikz}
\usepackage{braket}
\newtheorem{lemma}{Lemma}
\title{Perturbative anomalies in quantum mechanics}
\author{Maxim Gritskov, Andrey Losev, Saveliy Timchenko}
\date{Source: arXiv:2602.20968v5 (CC BY 4.0)}
\begin{document}
\maketitle

\noindent\emph{Source and licence.} Perturbative anomalies in quantum mechanics. Version arXiv:2602.20968v5 (CC BY 4.0), distributed by arXiv under the Creative Commons Attribution 4.0 International licence, http://creativecommons.org/licenses/by/4.0/, as recorded in the arXiv licence metadata for this exact version and shown on the abstract page https://arxiv.org/abs/2602.20968v5. Full licence notice: \texttt{licenses/arxiv-2602.20968-CC-BY-4.0.txt}.\par\smallskip

\noindent\textit{Editorial scope.} Counted unit: the article's lemma lemma:decomposition (source0.tex lines 239--274). The article's complete original proof of it is delivered with it (source0.tex lines 351--353, one \texttt{proof} environment of 3 lines, which the article places away from the statement). No mathematics is written, repaired or re-derived: the statement and the proof are the article's own lines, in the article's own notation, and no retained line is substituted or re-flowed.  Retained uncounted as the source-local context the proof needs: lines 224--227; lines 229--232; lines 275--326; lines 327--340.  Nothing was omitted from any retained span, so the package carries no omission record.  No retained line contains a citation, so the wrapper carries no bibliography and no bibliography entry.  No retained line contains a figure environment, an image-inclusion command or drawing code, so no image is delivered and the package declares no asset dependency.  Editorial formatting only, in addition: an AMS \texttt{amsart} preamble that restates the article's own declarations and macros used by the retained lines, namely the environments \texttt{lemma} on separate counters, together with the standard packages those lines need.  The retained environments print in this document as Lemma 1 in order of appearance, whereas the article prints the same items with the numbering of its own full document.  The complete native source of the article as distributed in the arXiv e-print for this exact version is bundled unchanged as \texttt{sources/195232/source0.tex}.
\begin{equation}
\label{CE}
    0\rightarrow \mathfrak{u}(V)\xlongrightarrow{d}(c^H\oplus c^S)\otimes \mathfrak{u}(V) \xlongrightarrow{d} c^H c^S\otimes \mathfrak{u}(V)\rightarrow 0,
\end{equation}

\begin{equation}
\label{CE_diffab}
    d=c^{H}\cdot\mathrm{ad}_{H}+c^{S}\cdot\mathrm{ad}_{S}.
\end{equation}

\begin{lemma} \label{lemma:acyclic}
    Consider $\bigoplus_{(a,\alpha)\neq(b,\beta)}\mathcal{B}_{(a,\alpha),(b,\beta)}^\bullet$. This subcomplex is acyclic.
    \begin{proof}
        It suffices to show this componentwise. Consider $0$-cycle $\omega$ in $\mathcal{B}^{\bullet}_{(a,\alpha),(b,\beta)}$:
        \begin{eqnarray}
            \notag d\omega=c^{H}\cdot \left(\frac{\mathrm{i}\lambda_{ab}}{2}\cdot \Pi_{V_{(a,\alpha)}}\tilde{\omega}\,\Pi_{V_{(b,\beta)}}-\frac{\mathrm{i}\lambda_{ab}}{2}\cdot \Pi_{V_{(b,\beta)}}\tilde{\omega}\,\Pi_{V_{(a,\alpha)}}\right)+\\+\,c^{S}\cdot \left(\frac{\mathrm{i}\mu_{\alpha\beta}}{2}\cdot \Pi_{V_{(a,\alpha)}}\tilde{\omega}\,\Pi_{V_{(b,\beta)}}-\frac{\mathrm{i}\mu_{\alpha\beta}}{2}\cdot \Pi_{V_{(b,\beta)}}\tilde{\omega}\,\Pi_{V_{(a,\alpha)}}\right)=0.
        \end{eqnarray}
        However, $\lambda_{ab}$ and $\mu_{\alpha\beta}$ are not equal to zero simultaneously. Suppose that $\lambda_{ab}\neq 0$:
        \begin{equation}
            \Pi_{V_{(a,\alpha)}}\tilde{\omega}\,\Pi_{V_{(b,\beta)}}- \Pi_{V_{(b,\beta)}}\tilde{\omega}\,\Pi_{V_{(a,\alpha)}}=0.
        \end{equation}
        It follows that $\omega=0$ and, therefore,
        \begin{equation}
        H^{0}(\mathcal{B}^{\bullet}_{(a,\alpha),(b,\beta)},d)=0.
        \end{equation}

 Now, let us consider the $1$-cocycle $\omega=c^{H}\omega_{H}+c^{S}\omega_{S}$:
        \begin{equation}
            d\omega=c^{H}c^{S}[H,\omega_{S}]-c^{H}c^{S}[S,\omega_{H}]=0.
        \end{equation}
        Note that again $\lambda_{ab}$ and $\mu_{\alpha\beta}$ cannot both be zero at the same time. Therefore
        \begin{eqnarray}
            \notag\frac{\mathrm{i}\lambda_{ab}}{2}\cdot \Pi_{V_{(a,\alpha)}}\tilde{\omega}_{S}\,\Pi_{V_{(b,\beta)}}-\frac{\mathrm{i}\lambda_{ab}}{2}\cdot \Pi_{V_{(b,\beta)}}\tilde{\omega}_{S}\,\Pi_{V_{(a,\alpha)}}=\\=\,\frac{\mathrm{i}\mu_{\alpha\beta}}{2}\cdot \Pi_{V_{(a,\alpha)}}\tilde{\omega}_{H}\,\Pi_{V_{(b,\beta)}}-\frac{\mathrm{i}\mu_{\alpha\beta}}{2}\cdot \Pi_{V_{(b,\beta)}}\tilde{\omega}_{H}\,\Pi_{V_{(a,\alpha)}}.
        \end{eqnarray}
        Suppose that $\lambda_{ab}\neq 0$. Then it follows from this equation that 
        \begin{equation}
            \omega_{S}=\frac{\mu_{\alpha\beta}}{\lambda_{ab}}\cdot \omega_{H}.
        \end{equation}
        Then we should find such a $0$-chain $\Omega$ that
        \begin{eqnarray}
            \notag d\Omega=\omega=c^{H}\omega_{H}+\frac{\mu_{\alpha\beta}}{\lambda_{ab}}\cdot c^{S}\omega_{H}=\\\notag=c^{H}\cdot \left(\frac{1}{2}\cdot \Pi_{V_{(a,\alpha)}}\tilde{\omega}_{H}\,\Pi_{V_{(b,\beta)}}+\frac{1}{2}\cdot \Pi_{V_{(b,\beta)}}\tilde{\omega}_{H}\,\Pi_{V_{(a,\alpha)}}\right)+\\ +\,\frac{\mu_{\alpha\beta}}{\lambda_{ab}}\cdot c^{S}\cdot \left(\frac{1}{2}\cdot \Pi_{V_{(a,\alpha)}}\tilde{\omega}_{H}\,\Pi_{V_{(b,\beta)}}+\frac{1}{2}\cdot \Pi_{V_{(b,\beta)}}\tilde{\omega}_{H}\,\Pi_{V_{(a,\alpha)}}\right).
        \end{eqnarray}
        Then $\Omega$ can be chosen as follows:
        \begin{equation}
            \Omega = \frac{1}{2\mathrm{i}\lambda_{ab}}\cdot \Pi_{V_{(a,\alpha)}}\tilde{\omega}_{H}\,\Pi_{V_{(b,\beta)}}-\frac{1}{2\mathrm{i}\lambda_{ab}}\cdot \Pi_{V_{(b,\beta)}}\tilde{\omega}_{H}\,\Pi_{V_{(a,\alpha)}}.
        \end{equation}
        The fact that $d\Omega=\omega$ can be verified by direct calculation. The fact that the cochain $\Omega\in\mathcal{B}_{(a,\alpha),(b,\beta)}$ was shown in \eqref{preimage}. Then we obtained that
        \begin{equation}
        H^{1}(\mathcal{B}^{\bullet}_{(a,\alpha),(b,\beta)},d)=0.
        \end{equation}

        Finally, it remains to show that every $2$-cochain $\omega$ has a preimage. To do this, we will again use the fact that $\lambda_{ab}$ and $\mu_{\alpha\beta}$ are not equal to zero at the same time. Let $\lambda_{ab}\neq 0$ again, then as a $1$-cochain, we will consider
        \begin{equation}
            \Omega=c^{S}\Omega_{S}=c^{S}\cdot\left(\frac{1}{2\mathrm{i}\lambda_{ab}}\cdot \Pi_{V_{(a,\alpha)}}\tilde{\omega}\,\Pi_{V_{(b,\beta)}}-\frac{1}{2\mathrm{i}\lambda_{ab}}\cdot \Pi_{V_{(b,\beta)}}\tilde{\omega}\,\Pi_{V_{(a,\alpha)}}\right).
        \end{equation}
        Then $d\Omega=c^{H}c^{S}[H,\Omega_{S}]$ and using the fact that $\Omega_{S}\in\mathcal{B}_{(a,\alpha),(b,\beta)}$ we obtain
        \begin{equation}
            d\Omega=c^{H}c^{S}\cdot\left(\frac{1}{2}\cdot \Pi_{V_{(a,\alpha)}}\tilde{\omega}\,\Pi_{V_{(b,\beta)}}+\frac{1}{2}\cdot \Pi_{V_{(b,\beta)}}\tilde{\omega}\,\Pi_{V_{(a,\alpha)}}\right)=\omega.
        \end{equation}
        It finally follows from this that $H^{2}(\mathcal{B}^{\bullet}_{(a,\alpha),(b,\beta)},d)=0$.
        \end{proof}
\end{lemma}

\begin{lemma} \label{prop:zero-diff}
    Let $\mathcal{Z}^{\bullet}=\bigoplus_{(a,\alpha)}\mathcal{B}_{(a,\alpha),(a,\alpha)}^{\bullet}$. Then $d|_{\mathcal{Z}^{\bullet}}=0$.
    \begin{proof}
        This immediately follows from the fact that for every $\mathcal{B}_{(a,\alpha),(a,\alpha)}$
        \begin{equation}
            \lambda_{aa}=\mu_{\alpha\alpha}=0.
        \end{equation}
        Therefore, the complex becomes a sequence of vector spaces with zero maps:
        \begin{equation}
            0\rightarrow\mathcal{Z}\xlongrightarrow{0}(c^H\oplus c^S)\otimes\mathcal{Z}\xlongrightarrow{0}(c^H c^S)\otimes\mathcal{Z}\rightarrow 0.
        \end{equation}
        The cohomology groups are simply the vector spaces themselves.
    \end{proof}
\end{lemma}

\begin{lemma} \label{lemma:decomposition}
    The Lie algebra $\mathfrak{u}(V)$ decomposes into a direct sum of vector spaces
    \begin{equation}
    \label{grading}
        \mathfrak{u}(V)=\bigoplus_{(a,\alpha),(b,\beta)}\mathcal{B}_{(a,\alpha),(b,\beta)}
    \end{equation}
    where $\mathcal{B}_{(a,\alpha),(b,\beta)}$ is the image of the linear operator $\pi_{(a,\alpha)}^{(b,\beta)}$:
    \begin{equation}
        \pi_{(a,\alpha)}^{(b,\beta)}(x)=\frac{1}{2}\cdot \Pi_{V_{(a,\alpha)}}x\,\Pi_{V_{(b,\beta)}}+\frac{1}{2}\cdot \Pi_{V_{(b,\beta)}}x\,\Pi_{V_{(a,\alpha)}}.
    \end{equation}
    Then, the complex $\eqref{CE}$ decomposes into a direct sum of the complexes $\mathcal{B}_{(a,\alpha),(b,\beta)}^{\bullet}$:
    \begin{equation}
    \label{CE_decomp}
        0\rightarrow \mathcal{B}_{(a,\alpha),(b,\beta)}\xlongrightarrow{d}(c^H\oplus c^S)\otimes \mathcal{B}_{(a,\alpha),(b,\beta)} \xlongrightarrow{d} c^H c^S\otimes \mathcal{B}_{(a,\alpha),(b,\beta)}\rightarrow 0.
    \end{equation}
    \begin{proof}
    Consider the action of the differential \eqref{CE_diffab} on the general element of \eqref{CE_decomp}:
        \begin{eqnarray}
            \notag d(w+c^{H}x+c^{S}y+c^{H}c^{S}z)=[c^H H+c^S S, w]+\{c^H H+c^S S,c^{H}x+c^{S}y\}=\\=c^{H}\cdot [H,w]+c^{S}\cdot [S,w]+c^{H}c^{S}\cdot[H,y]-c^{H}c^{S}\cdot[S,x]\,.
        \end{eqnarray}
    Then it is sufficient to show that for any element $\omega\in\mathcal{B}_{(a,\alpha),(b,\beta)}$, it is true that 
        \begin{equation}
            [H,\omega]\in\mathcal{B}_{(a,\alpha),(b,\beta)},\, [S,\omega]\in\mathcal{B}_{(a,\alpha),(b,\beta)}.
        \end{equation}
    We will prove this for $[H,\omega]$. Let there exist such an $\tilde{\omega}$ that $\omega=\pi_{(a,\alpha)}^{(b,\beta)}(\tilde{\omega})$, then
        \begin{equation}
            [H,\omega]=\frac{\mathrm{i}\lambda_{ab}}{2}\cdot \Pi_{V_{(a,\alpha)}}\tilde{\omega}\,\Pi_{V_{(b,\beta)}}-\frac{\mathrm{i}\lambda_{ab}}{2}\cdot \Pi_{V_{(b,\beta)}}\tilde{\omega}\,\Pi_{V_{(a,\alpha)}}\in \mathfrak{u}(V),
        \end{equation}
    where $\lambda_{ab}=\lambda_{a}-\lambda_{b}$. But then it is easy to find its $\pi_{(a,\alpha)}^{(b,\beta)}$-preimage:
        \begin{equation}
        \label{preimage}
            [H,\omega]=\pi_{(a,\alpha)}^{(b,\beta)}(\mathrm{i}\lambda_{ab}\cdot \Pi_{V_{(a,\alpha)}}\tilde{\omega}\,\Pi_{V_{(b,\beta)}}-\mathrm{i}\lambda_{ab}\cdot\Pi_{V_{(b,\beta)}}\tilde{\omega}\,\Pi_{V_{(a,\alpha)}}).
        \end{equation}
        Then it follows that $[H,\omega]\in\mathcal{B}_{(a,\alpha),(b,\beta)}$. For $[S,\omega]$, the proof is similar.
    \end{proof}
\end{lemma}

    \begin{proof}
        This immediately follows from the lemmas~\ref{lemma:decomposition}, ~\ref{lemma:acyclic} and ~\ref{prop:zero-diff}.
    \end{proof}

\end{document}
